\documentclass[twocolumn,pra,amsmath,amssymb,superscriptaddress]{revtex4-2}

\usepackage{graphicx}
\usepackage{hyperref}
\usepackage{booktabs}
\usepackage{xcolor}
\usepackage{amsthm}

\newtheorem{theorem}{Theorem}

\begin{document}

\title{The Universal/Coulomb Partition of the GeoVac Spectral Triple\\
{\large What belongs to angular momentum, what belongs to $-Z/r$,\\
and why the difference matters}}

\author{J.~Loutey}
\affiliation{Independent Researcher, Kent, Washington}

\date{May 2, 2026}

\begin{abstract}
The GeoVac framework~\cite{paper7,paper14} is a discrete spectral
triple in the sense of Connes~\cite{connes_book1994} and the
graph-spectral-triple lineage of Marcolli and van
Suijlekom~\cite{marcolli_vs2014}: an algebra $\mathcal{A}$ of
functions on the Fock-projected discrete graph, a Hilbert space
$\mathcal{H}$ of scalar and spinor states, and a Dirac operator $D$
inherited from the Camporesi--Higuchi spectrum on $S^3$.  This paper
identifies a clean structural decomposition of the framework's
results into two disjoint sectors:
\begin{itemize}
\item A \emph{universal sector}, depending only on $\mathcal{A}$ and
  the angular structure of the basis.  The angular sparsity theorem
  of Paper~22~\cite{paper22} --- ERI density $1.44\%$ at $l_{\max}=3$,
  bit-identical across Coulomb, harmonic oscillator, Woods--Saxon,
  square well, and Yukawa potentials --- is the canonical example.
\item A \emph{potential-specific sector}, depending on the choice
  of $D$.  The Fock projection rigidity theorem of
  Paper~23~\cite{paper23} ($S^3$ projection unique to $-Z/r$ by
  $\mathrm{SO}(4)$ symmetry) and the Bargmann--Segal rigidity
  theorem of Paper~24~\cite{paper24} (the 3D isotropic harmonic
  oscillator is the unique central potential whose spectrum arises
  from the Euler operator on $H^2(S^5)$ restricted to symmetric
  $\mathrm{SU}(3)$ irreps) jointly establish the partition's
  sharpness.
\end{itemize}
The two rigidity theorems sit on opposite sides of an operator-order
divide: Coulomb requires a second-order Riemannian Laplacian and a
nonlinear stereographic projection that introduces calibration $\pi$;
the HO uses a first-order complex Euler operator and a linear-affine
Bargmann projection that is bit-exactly $\pi$-free.  This asymmetry is
not an accident of two examples; it is a structural feature of
sphere-quotient discretizations.  The Sprint 5 result on $S^5$ gauge
structure~\cite{paper25} sharpens the Coulomb/HO asymmetry from two
layers to three (spectrum-computing role of $L_0$, calibration $\pi$,
and Wilson physical content), all three stratifying along the same
operator-order divide.  We frame the partition in spectral-triple
language and identify which features of GeoVac transfer to non-Coulomb
systems for free, which require system-specific rederivation, and what
this implies for nuclear, oscillator, and gauge-theoretic
extensions of the framework.
\end{abstract}

\maketitle

%=================================================================
\section{Introduction}
\label{sec:intro}
%=================================================================

The GeoVac framework~\cite{paper0,paper1,paper7} discretizes
quantum-mechanical Hamiltonians on natural geometries by replacing
the continuous Laplace--Beltrami operator with the graph Laplacian
on a packed lattice of quantum-number labels.  For the hydrogen atom,
the result is exact: the discrete spectrum $\lambda_n = -(n^2 - 1)$
recovers the hydrogenic Rydberg series under Fock's stereographic
projection $p_0 = \sqrt{-2E}$ from $S^3$ to momentum
space~\cite{fock1935,bander_itzykson1966}.  The framework has been
extended to multi-electron atoms~\cite{paper13},
diatomics~\cite{paper11,paper15}, polyatomics~\cite{paper17},
nuclear shell model Hamiltonians~\cite{paper23}, and the
three-dimensional isotropic harmonic oscillator~\cite{paper24}.  In
each case the same combinatorial machinery --- Wigner $3j$ symbols,
Gaunt selection rules, $jj$-coupled angular reduction --- produces a
sparse one- and two-body Hamiltonian with the same structural
features.  Paper~22's central observation~\cite{paper22} is that this
sparsity is not particular to any of these systems: it is a property
of the angular momentum algebra and persists across all spherical
fermion bases, regardless of the radial potential.

The same machinery, however, does \emph{not} produce identical
transcendental content.  The Coulomb framework on $S^3$ is honestly
$\pi$-laden at every level above the bare graph: the kinetic
calibration $\kappa = -1/16$ ratio is set by the Fock projection
Jacobian, the $\alpha$-conjecture combination $K = \pi(B + F -
\Delta)$~\cite{paper2_alpha} carries an explicit factor of $\pi$ as
the Hopf-bundle measure $\mathrm{Vol}(S^2)/4 = \pi$, and the spectral
zeta values appear at even integer arguments via the Jacobi-theta
inversion formula~\cite{paper18,paper28}.  The 3D harmonic oscillator
discretized via Bargmann--Segal on $S^5$ is by contrast bit-exactly
$\pi$-free in exact rational arithmetic at every finite
$N_{\max}$~\cite{paper24}: every adjacency, every dipole transition
amplitude, every spectral eigenvalue is a rational number or finite
algebraic extension.  This is not a numerical accident; Paper~24
establishes that no $\pi$ enters the discrete graph because no
calibration step has been performed.  The harmonic oscillator's
Bargmann transform is unitary by construction, requires no
stereographic distortion, and produces a linear spectrum
$\hbar\omega(N + 3/2)$ with no transcendental coefficient anywhere.

These two observations --- universal angular sparsity and
potential-specific transcendental content --- point at a structural
decomposition of the GeoVac framework that, once stated explicitly,
clarifies what the framework does and does not transfer across
physical systems.  The decomposition is the standard separation of a
spectral triple $(\mathcal{A}, \mathcal{H}, D)$ into its algebra
$\mathcal{A}$ and its Dirac operator $D$, in the sense of
Connes~\cite{connes_book1994} and the graph-spectral-triple framework
of Marcolli and van Suijlekom~\cite{marcolli_vs2014}.  In this
language, angular sparsity belongs to $\mathcal{A}$ (it is a
selection rule on which products of basis functions are nonzero,
independent of any operator); rigidity, calibration $\pi$, and the
particular spectrum belong to $D$ (they encode the choice of
potential and projection).

The contribution of this paper is to state the partition explicitly,
give it a name, and trace its consequences.  Its scientific content
is small: every theorem cited is established in a previous paper of
this series.  Its organizational content is large: the partition
explains \emph{which} GeoVac results are universal across spherical
fermion problems and \emph{which} are particular to the Coulomb
potential, and it provides a checklist for extending the framework
to new potentials without re-deriving universal results.

The paper is structured as follows.
Section~\ref{sec:setup} introduces the GeoVac framework as a discrete
spectral triple in the Marcolli--van Suijlekom sense.
Section~\ref{sec:universal} restates Paper~22's angular sparsity
theorem as the canonical content of the universal sector.
Section~\ref{sec:specific} reviews the two known potential-specific
rigidity theorems (Fock and Bargmann--Segal) and shows that they sit
on opposite sides of an operator-order divide.
Section~\ref{sec:transcendental} traces the calibration-$\pi$
asymmetry to the operator-order distinction and the projection type
(nonlinear vs.\ linear-affine).
Section~\ref{sec:three-layers} reports the Sprint~5 result that the
Coulomb/HO asymmetry has three layers, not two, and gives the
predictive checklist that follows.
Section~\ref{sec:implications} states what features of GeoVac
transfer to non-Coulomb systems for free and what does not.
Section~\ref{sec:open} closes with open questions about the
extension to non-central potentials and gauge sectors.

\medskip\noindent\textbf{Scope of contribution.}  This paper is not a
new theorem.  It is a structural framing of existing
results~\cite{paper22,paper23,paper24,paper25} in spectral-triple
language.  No new computations are performed.  The framing is
intended to clarify which existing results extend without
modification to non-Coulomb systems and which require explicit
recomputation, and to place the GeoVac framework in the
graph-spectral-triple lineage of Connes, Marcolli, and van
Suijlekom~\cite{connes_book1994,marcolli_vs2014,perez_sanchez2024}.
The dual-description rule of the framework --- in particular that the
identification of the discrete graph with the continuum
Laplace--Beltrami operator is a proven mathematical equivalence and
not an ontological claim about which description is more
fundamental --- is preserved throughout.

%=================================================================
\section{Setup: GeoVac as a Discrete Spectral Triple}
\label{sec:setup}
%=================================================================

A spectral triple in the sense of Connes~\cite{connes_book1994} is a
data triple
\[
  (\mathcal{A}, \mathcal{H}, D)
\]
where $\mathcal{A}$ is a $*$-algebra represented as bounded
operators on a Hilbert space $\mathcal{H}$, and $D$ is a self-adjoint
operator on $\mathcal{H}$ with compact resolvent and bounded
commutators $[D, a]$ for all $a \in \mathcal{A}$.  In the commutative
case $\mathcal{A} = C^\infty(M)$ for a closed Riemannian spin manifold
$M$ and $D$ is the Dirac operator on $M$, the triple recovers the
geometry of $M$ via Connes' reconstruction theorem.  In the
noncommutative case, the triple is treated as defining the geometry,
and the spectral action principle of Chamseddine and
Connes~\cite{chamseddine_connes1997} extracts physics by computing
$\mathrm{Tr}\, f(D^2/\Lambda^2)$ for suitable test functions~$f$ and
cutoffs~$\Lambda$.

Marcolli and van Suijlekom~\cite{marcolli_vs2014} introduced
\emph{gauge networks}: a discrete variant of the spectral triple
where $\mathcal{A}$ is supported on the vertices of a graph,
$\mathcal{H}$ is a finite-dimensional vector space (one fiber per
vertex), and the gauge connection lives on the edges.  Their main
result is that the Chamseddine--Connes spectral action on a gauge
network coincides with the Wilson lattice gauge action of the
underlying graph; the Perez-Sanchez correction~\cite{perez_sanchez2024}
clarifies that the continuum limit recovers Yang--Mills without
Higgs.

The GeoVac framework fits this lineage almost verbatim.  Specifically:

\begin{description}
\item[$\mathcal{A}$:]  the $*$-algebra of functions on the discrete
  Fock-projected graph, with vertices labeled by quantum numbers
  $(n, l, m)$ for the scalar sector and $(n, \kappa, m_j)$ for the
  spinor sector~\cite{paper7}.  Multiplication is pointwise; the
  involution is complex conjugation.  $\mathcal{A}$ acts on
  $\mathcal{H}$ by diagonal multiplication operators
  $a_v \cdot \psi_v = (a \cdot \psi)_v$.
\item[$\mathcal{H}$:]  the finite-dimensional Hilbert space of
  scalar (or spinor) wavefunctions on the graph.  Decomposes as
  $\mathcal{H} = \mathcal{H}^{\mathrm{scalar}} \oplus
  \mathcal{H}^{\mathrm{spinor}}$ with the spinor sector supplied by
  the Camporesi--Higuchi construction~\cite{paper28}.
\item[$D$:]  the Dirac operator on $\mathcal{H}^{\mathrm{spinor}}$
  with eigenvalues $|\lambda_n^{\mathrm{D}}| = n + 3/2$ and
  degeneracies $g_n^{\mathrm{D}} = 2(n+1)(n+2)$, restricted to a
  finite cutoff $n \le n_{\max}$~\cite{paper28}.  The squared Dirac
  operator $D^2$ acts on the scalar sector with the
  Lichnerowicz formula~\cite{lichnerowicz1963}
  \[
    D^2 = \Delta_{\mathrm{LB}} + \tfrac{R_{\mathrm{scalar}}}{4}
    = \Delta_{\mathrm{LB}} + \tfrac{3}{2}
  \]
  on the unit $S^3$, recovering the eigenvalues
  $|\lambda_n|^2 = (n + 3/2)^2 = n(n+2) + 9/4$.
\end{description}

The non-abelian gauge structure of Papers~25 and 30 enters as inner
derivations of an almost-commutative extension
$\mathcal{A} \otimes M_n(\mathbb{C})$, with
Paper~25~\cite{paper25} corresponding to the $n = 1$ (U(1)) maximal
torus and Paper~30~\cite{paper30} to the $n = 2$ (SU(2)) full
non-abelian slice.  The five action principles documented in
Paper~30~\S6 (entropy maximization on 0-cochains, Rayleigh--Ritz on
the graph Laplacian, Wilson action on the gauge sector, Ihara zeta
path integral, spectral action) each organize a distinct sector of the
spectral triple data.

\medskip\noindent\textbf{The decomposition.} Once the data
$(\mathcal{A}, \mathcal{H}, D)$ are fixed, every observable in the
framework is a function of these three pieces and a finite list of
auxiliary structures (the cutoff $n_{\max}$, the gauge group, the
test function $f$ for the spectral action).  The partition we develop
in this paper observes that the GeoVac results split cleanly along
the $\mathcal{A}$ versus $D$ line:
\begin{itemize}
\item Results that depend only on $\mathcal{A}$ (and the angular
  decomposition of its basis) are \emph{universal across central
  potentials}: they hold for any spectral triple sharing the same
  algebra and basis, regardless of which Dirac operator is chosen.
\item Results that depend on the specific form of $D$ (or
  equivalently, on the radial potential $V(r)$) are
  \emph{potential-specific}: they hold for the Coulomb problem (or
  the harmonic oscillator, or a Woods--Saxon nucleus) but not for
  arbitrary central potentials.
\end{itemize}
The decomposition is not novel as a piece of spectral triple
formalism --- it is the standard $\mathcal{A}$/$D$ split.  The
content of this paper is to identify which GeoVac theorems live on
each side and what this implies for the framework's extension to new
systems.

%=================================================================
\section{The Universal Sector: Angular Sparsity}
\label{sec:universal}
%=================================================================

The canonical content of the universal sector is the angular sparsity
theorem of Paper~22~\cite{paper22}.  We restate it in spectral-triple
language to make the partition framework explicit.

\subsection{Restatement of the angular sparsity theorem}

Let $(\mathcal{A}, \mathcal{H}, D)$ be a GeoVac spectral triple in
which $\mathcal{H}$ is decomposed in the spherical-harmonic basis
$|n, l, m\rangle$ and the rotationally invariant two-body operator
$\hat{V}_{12} = V(|\mathbf{r}_1 - \mathbf{r}_2|)$ admits the Legendre
expansion
\[
  V(|\mathbf{r}_1 - \mathbf{r}_2|)
  = \sum_{k \ge 0} v_k(r_1, r_2)\, P_k(\cos\theta_{12}).
\]
Then the two-body matrix elements factorize~\cite{paper22} as
\begin{widetext}
\[
  \langle n_1 l_1 m_1, n_2 l_2 m_2\,|\,\hat{V}_{12}\,|
  \,n_3 l_3 m_3, n_4 l_4 m_4 \rangle
  = \sum_k c^k(l_1 m_1; l_3 m_3)\, c^k(l_2 m_2; l_4 m_4)\,
  R^k(n_1 l_1, n_2 l_2; n_3 l_3, n_4 l_4),
\]
\end{widetext}
where the Gaunt coefficients $c^k$ are pure functions of $(l, m)$
quantum numbers and the radial integrals $R^k$ depend on $V$.  The
zero structure of the matrix-element tensor is determined by the
Gaunt selection rules alone:
\begin{itemize}
\item parity: $l_1 + l_3 + k$ and $l_2 + l_4 + k$ both even,
\item triangle inequalities $|l_1 - l_3| \le k \le l_1 + l_3$ and
  $|l_2 - l_4| \le k \le l_2 + l_4$,
\item magnetic conservation $m_1 + m_2 = m_3 + m_4$.
\end{itemize}
None of these conditions reference $V$, $R^k$, or any radial quantum
number.  Theorem~2 of Paper~22 then states that the ERI density
$D(l_{\max})$ is a function only of the angular momentum cutoff:
\begin{equation}
  D(l_{\max}) =
  \begin{cases}
    7.81\% & l_{\max} = 1 \\
    2.76\% & l_{\max} = 2 \\
    1.44\% & l_{\max} = 3 \\
    0.90\% & l_{\max} = 4 \\
    0.62\% & l_{\max} = 5
  \end{cases}
  \label{eq:eri_density}
\end{equation}
Paper~22 verified this numerically across five potentials (Coulomb,
harmonic oscillator, Woods--Saxon, square well, Yukawa) at
$l_{\max} = 2$, finding bit-identical zero patterns and
$D = 2.76\%$ in every case.

\subsection{Spectral-triple reading}

In the partition framework, the angular sparsity theorem is a
statement about the algebra $\mathcal{A}$ and the $\mathrm{SO}(3)$
representation structure of its basis decomposition.  The set of
nonzero matrix elements
\[
  \mathcal{N}(\mathcal{A}, l_{\max})
  = \{(a, b, c, d) \in \mathcal{A}^{\otimes 4} :
     \langle a b | \hat{V}_{12} | c d \rangle \ne 0\}
\]
is determined by the Gaunt triangle and parity conditions, which
depend only on the basis labels $(l, m)$ and not on $D$.  Equivalently,
$\mathcal{N}$ is a function of $\mathcal{A}$ alone, in the precise
sense that any change of $D$ (i.e., any change of radial potential)
that preserves the spherical-harmonic basis decomposition of
$\mathcal{H}$ leaves $\mathcal{N}$ unchanged.

This is the formal sense in which angular sparsity is universal.  The
spinor extension developed in Paper~22's spinor-block angular sparsity
section~\cite{paper22} carries the same property: the $jj$-coupled angular reduction~\cite{Dyall} preserves
the radial-angular factorization and the spinor ERI density
$d_{\mathrm{spinor}}(l_{\max})$ depends only on $l_{\max}$, not on the
choice of $D$.  Thus the universal sector includes both scalar and
spinor sparsity bounds.

\subsection{What the universal sector contains}

The universal sector encompasses every result of the GeoVac framework
that can be derived from the algebra $\mathcal{A}$, the angular
decomposition of $\mathcal{H}$, and the Wigner $3j$ algebra,
without invoking the specific spectrum of $D$.  This includes:

\begin{itemize}
\item Angular sparsity bounds: Eq.~\eqref{eq:eri_density} above.
\item The composed-block $O(Q^{2.5})$ Pauli scaling of
  Paper~14~\cite{paper14}, which depends on the block-diagonal
  structure of two-body matrix elements in a composed orbital basis
  and the Gaunt selection rules.  Verified across all 28 molecules
  in the GeoVac library.
\item Selection rules for transitions, oscillator strengths, and
  multipole moments (the Wigner--Eckart content of any operator
  expanded in spherical tensor components).
\item Block-diagonal structure of two-body interaction tensors in
  the composed architecture~\cite{paper17}.
\item The $jj$-coupled qubit encoding sparsity for relativistic
  composed Hamiltonians~\cite{paper14}~\S V (Tier~2 spinor
  extension).
\end{itemize}

Each of these results was originally derived in the Coulomb context,
but the derivation never used a Coulomb-specific feature of $D$
beyond the spherical-harmonic basis decomposition.  They transfer to
any new central potential without re-derivation, contingent only on
the new system having the same basis.

\medskip\noindent\textbf{Practical consequence.}  When extending the
GeoVac framework to a new physical system (nuclear shell model,
quantum dot, oscillator basis, screened Coulomb, etc.), the universal
sector results may be \emph{quoted} from Paper~22 and Paper~14
without recomputation.  Verifying the new system gives the same
sparsity is then a single sanity check, not a research program.

%=================================================================
\section{The Potential-Specific Sector: Two Rigidity Theorems}
\label{sec:specific}
%=================================================================

The potential-specific sector contains results that depend on the
explicit form of $D$.  The two most consequential examples in the
GeoVac framework to date are the Fock projection rigidity theorem
(Paper~23) and the Bargmann--Segal rigidity theorem (Paper~24).
Both are uniqueness theorems: each identifies a single central
potential as the unique system whose spectrum arises from a specific
sphere-quotient construction.  The two theorems sit on opposite
sides of an operator-order divide that we make explicit in
Section~\ref{sec:transcendental}.

\subsection{Fock projection rigidity}
\label{sec:fock_rigidity}

Paper~23~\cite{paper23} states the following:

\begin{theorem}[Fock projection rigidity, Paper~23]
\label{thm:fock_rigidity}
The $S^3$ Fock projection $p_0 = \sqrt{-2 E_n}$ maps a one-particle
central-field Hamiltonian onto the free Laplacian on $S^3$ if and
only if the spectrum $E_{n\ell}$ is independent of $\ell$ within
each $n$.  The Coulomb potential $-Z/r$ is the unique central
potential exhibiting this accidental degeneracy, as a consequence of
its hidden $\mathrm{SO}(4)$ symmetry (the Runge--Lenz vector).  Any
deformation $V(r) = -Z/r + \delta V(r)$ with $\delta V \not\equiv 0$
lifts the $\ell$-degeneracy, breaking the conformal equivalence
between the radial Schr\"odinger problem and the $S^3$ graph
Laplacian.
\end{theorem}

The proof~\cite{paper23} amounts to observing that Fock's
stereographic construction in momentum space requires the hydrogenic
relation $E_n \propto -1/n^2$ with no $\ell$-dependence, which fails
for any potential without the $\mathrm{SO}(4)$ symmetry of the pure
Coulomb problem~\cite{bander_itzykson1966}.

\medskip\noindent\textbf{What rigidity rules out.}  The theorem
forbids the following extensions of GeoVac to non-Coulomb central
potentials:
\begin{itemize}
\item Fock projection $p_0 = \sqrt{-2E}$ does not map the radial
  Schr\"odinger equation onto the free Laplacian on $S^3$.
\item The integer eigenvalue spectrum $\lambda_n = -(n^2 - 1)$ of
  the graph Laplacian does not coincide with the rescaled energy
  spectrum.
\item The $\kappa = -1/16$ kinetic constant has no analog in the
  new system.
\end{itemize}
What rigidity does \emph{not} rule out is the angular machinery: the
Wigner $3j$ algebra, the Gaunt integrals, the block structure of the
two-body interaction tensor.  This is the universal-sector content
of Paper~22, and it transfers to any central potential.

\subsection{Bargmann--Segal rigidity}
\label{sec:bs_rigidity}

Paper~24~\cite{paper24} establishes the structural dual of
Theorem~\ref{thm:fock_rigidity}.

\begin{theorem}[Bargmann--Segal rigidity for the 3D HO, Paper~24]
\label{thm:bs_rigidity}
The three-dimensional isotropic harmonic oscillator
$V(r) = \tfrac{1}{2} m\omega^2 r^2$ is the unique central-field
potential whose bound-state spectrum arises from the Euler operator
on the Hardy space $H^2(S^5)$ restricted to the symmetric
irreducible representations $(N, 0)$ of $\mathrm{SU}(3)$.
\end{theorem}

The proof has three classical inputs~\cite{paper24}:
Jauch--Hill~\cite{jauchhill1940} (the 3D isotropic HO is the unique
central potential with $\mathrm{SU}(3)$ dynamical symmetry);
Borel--Weil~\cite{borelweil} (the symmetric $(N, 0)$ irreps coincide
with the homogeneous degree-$N$ subspaces of $H^2(S^5)$); and
Bargmann--Segal--Hall~\cite{bargmann1961,hall1994} (the unique
unitary equivalence between $L^2(\mathbb{R}^3)$ and the
Bargmann--Fock space, mapping the HO Hamiltonian to
$\hbar\omega(\hat{N} + 3/2)$).

\medskip\noindent\textbf{The structural dual.}  The two rigidity
theorems describe analogous but distinct constructions.  We tabulate
the parallels in Table~\ref{tab:rigidity_dual}.

\begin{table}[h]
\caption{The two known central-potential rigidity theorems in the
GeoVac framework.  Each identifies a unique central potential
matching a specific sphere-quotient construction.}
\label{tab:rigidity_dual}
\begin{ruledtabular}
\begin{tabular}{lll}
Feature & Coulomb (Paper~23) & 3D HO (Paper~24) \\
\hline
Sphere & $S^3$ & $S^5$ \\
Geometry & Riemannian & Complex (Hardy) \\
Symmetry & $\mathrm{SO}(4)$ & $\mathrm{SU}(3)$ \\
Operator & 2nd-order Laplacian & 1st-order Euler \\
Spectrum & $-1/n^2$ (quadratic) & $N + 3/2$ (linear) \\
Projection & $p_0 = \sqrt{-2E}$ (nonlinear) & linear-affine \\
Calibration & $\kappa = -1/16$, $\pi$-laden & $\pi$-free \\
\end{tabular}
\end{ruledtabular}
\end{table}

The two theorems are not connected by any direct mathematical bridge
(there is no central potential that satisfies both simultaneously).
What they share is a structural form: each pairs a unique potential
with a unique projection construction, and each forbids any
modification of the potential without breaking the construction.

\subsection{The shared structure: rigidity as $D$-dependence}

Both rigidity theorems are statements about $D$.  They say nothing
about $\mathcal{A}$: the algebra of functions on the graph, the
spherical-harmonic basis decomposition, the Gaunt selection rules
remain valid in either case.  What they say is that the specific form
of $D$ (the spectrum, the projection used to extract physical
energies, the calibration constants) is uniquely determined by the
choice of central potential.  Two different central potentials
produce two different spectral triples that share an algebra but
differ in $D$.

This is the formal sense in which rigidity belongs to the
potential-specific sector.  Any framework result that depends on
the explicit form of $D$ --- the integer spectrum, the calibration
$\kappa = -1/16$, the $\alpha$-conjecture combination
$K = \pi(B + F - \Delta)$~\cite{paper2_alpha} --- is bound to one
specific central potential and does not transfer.

%=================================================================
\section{The Transcendental Asymmetry}
\label{sec:transcendental}
%=================================================================

The two rigidity theorems differ in more than the choice of sphere.
Coulomb is honestly $\pi$-laden at every level above the bare graph;
the harmonic oscillator is bit-exactly $\pi$-free at every level.
This asymmetry is not an accident of two examples; it traces directly
to the operator-order divide in Table~\ref{tab:rigidity_dual}.

\subsection{Operator order and projection type}

The Coulomb construction uses the second-order Laplace--Beltrami
operator $\Delta_{S^3}$ on the unit three-sphere, with eigenvalues
$\lambda_n = n(n+2)$ (quadratic in $n$).  The Fock projection
$p_0 = \sqrt{-2 E_n}$ extracts the physical energy $E_n = -1/(2n^2)$
by matching $p_0$ to the radial scale of momentum-space wavefunctions
on $S^3$.  This projection is \emph{nonlinear} in the eigenvalue:
\[
  E_n = -\tfrac{1}{2 n^2}
  \;\;\Longrightarrow\;\;
  E_n^{-1/2} = \sqrt{-2 E_n} \cdot n
  = p_0(n) \cdot n,
\]
and the inverse relation $n \mapsto E_n$ involves an inverse square
root.  The calibration step from graph eigenvalue to physical energy
introduces a factor that traces to the Riemannian volume of $S^3$
(via the Lichnerowicz constant on $S^3$ and the Hopf-bundle measure
$\mathrm{Vol}(S^2)/4 = \pi$); this is the source of the $\pi$ that
appears in $\kappa = -1/16$ as $1/(16\pi^0)$ (specifically, the
ratio $\sqrt{\pi}/\sqrt{\pi}$ cancels but the factor of $4$ is set
by the projection Jacobian)~\cite{paper7,paper24}.

The harmonic oscillator construction uses the first-order Euler
operator $\hat{N}$ on the Hardy space $H^2(S^5)$, with eigenvalues
$N$ (linear in the holomorphic degree).  The
Bargmann--Segal--Hall projection is unitary and \emph{linear-affine}
in the eigenvalue:
\[
  E_N = \hbar\omega(N + 3/2),
\]
and the inverse relation $N \mapsto E_N$ involves no square roots,
no rescalings, and no Riemannian volume factors.  The calibration
step requires no $\pi$ because no Riemannian volume is integrated
to extract the physical scale.  The single $\pi$ in the
Bargmann--Segal ecosystem is the Gaussian normalization $\pi^{-3}$
of the continuous Bargmann measure, which the discrete graph never
invokes.

\subsection{The general structural rule}

The asymmetry generalizes.  The presence or absence of calibration
$\pi$ in a sphere-quotient discretization correlates with the
operator order and projection type as follows:

\begin{description}
\item[Second-order Riemannian, nonlinear projection.]  Calibration
  $\pi$ \emph{is} structurally forced.  The projection step
  integrates against a Riemannian volume form
  $d\mathrm{Vol}_M = \sqrt{|g|}\, d^n x$, which contributes
  $\mathrm{Vol}(M)$ to overlap normalizations and $\pi$ to spectral
  determinants via the Jacobi-theta inversion formula~\cite{paper28}.
  Even-power $\pi^{2k}$ enters via spectral zeta values
  $\zeta_R(2k)$, which are rational multiples of $\pi^{2k}$.
\item[First-order complex, linear-affine projection.]  Calibration
  $\pi$ is \emph{not} forced.  The projection is unitary and the
  spectrum is exact-linear in the holomorphic degree; no Riemannian
  volume is integrated and no square roots are taken.  Discrete
  graph quantities are expressible as exact rationals or finite
  algebraic extensions of $\mathbb{Q}$~\cite{paper24}.
\end{description}

This is the operator-order axis of the Paper~18 transcendental
taxonomy~\cite{paper18}.  The QED-on-$S^3$ work
(Paper~28~\cite{paper28}) confirmed it via the T9 theorem: the
spectral zeta function of the squared Dirac operator $D^2$ is a
two-term polynomial in $\pi^2$ with rational coefficients at every
integer $s$ (operator order 2 $\to$ $\pi^{\mathrm{even}}$).  The
parity discriminant theorem of Paper~28 confirms the dual: odd-order
operators produce odd-zeta content $\zeta_R(3), \zeta_R(5), \ldots$
via half-integer Hurwitz reduction.

\subsection{Where calibration $\pi$ does and does not appear}

The structural rule predicts which framework results carry $\pi$ and
which do not.  We tabulate examples in
Table~\ref{tab:pi_distribution}.

\begin{table}[h]
\caption{Where calibration $\pi$ appears in GeoVac results.  Items
above the divide depend on second-order Riemannian operators with
nonlinear projections; items below depend on first-order complex
operators with linear projections.  Universal-sector results (angular
sparsity, $jj$-coupled selection rules) are $\pi$-free regardless,
because they involve no projection.}
\label{tab:pi_distribution}
\begin{ruledtabular}
\begin{tabular}{ll}
Result & $\pi$? \\
\hline
\multicolumn{2}{l}{\emph{Universal sector (algebra only):}} \\
Angular sparsity~\cite{paper22} & no \\
Composed Pauli scaling~\cite{paper14} & no \\
Wigner $3j$ algebra & no \\
\hline
\multicolumn{2}{l}{\emph{Coulomb-specific (2nd-order, nonlinear proj.):}} \\
$\kappa = -1/16$~\cite{paper7} & yes \\
Spectral zeta $\zeta_{D^2}(s)$~\cite{paper28} & yes ($\pi^{2k}$) \\
Hopf-bundle measure $\mathrm{Vol}(S^2)/4$~\cite{paper2_alpha} & yes \\
$\alpha$-conjecture combination $K$~\cite{paper2_alpha} & yes \\
\hline
\multicolumn{2}{l}{\emph{HO-specific (1st-order, linear proj.):}} \\
Bargmann--Segal graph adjacency~\cite{paper24} & no \\
HO spectrum $\hbar\omega(N + 3/2)$ & no \\
Bargmann transition amplitudes & no \\
\end{tabular}
\end{ruledtabular}
\end{table}

\medskip\noindent\textbf{Implication for nuclear extensions.}
Paper~23's nuclear shell model qubit Hamiltonians use the harmonic
oscillator basis with the Minnesota nucleon--nucleon
potential~\cite{paper23}.  By the structural rule, these Hamiltonians
should be $\pi$-free in the angular and basis structure (HO is the
1st-order side); calibration $\pi$ enters only if the projection step
to physical energies invokes a Riemannian volume.  Paper~23's
construction works directly in the HO basis without a stereographic
projection, so the prediction is that the nuclear qubit Hamiltonians
are $\pi$-free.  Inspection of Paper~23 confirms this: every
coefficient is a rational combination of Minnesota Gaussians,
Moshinsky--Talmi brackets, and HO matrix elements; no transcendental
coefficient appears.

%=================================================================
\section{The Three-Layer Sharpening (Sprint~5)}
\label{sec:three-layers}
%=================================================================

The Sprint~5 investigation~\cite{paper25} of $S^5$ gauge structure
sharpened the Coulomb/HO asymmetry from two layers
(spectrum-computing and calibration $\pi$) to three layers, with the
third being the Wilson physical content of the gauge sector.  We
summarize the three-layer partition here as the predictive
checklist for extending the framework to any new central potential.

\subsection{The three layers}

\begin{description}
\item[Layer 1: Spectrum-computing role of the node Laplacian $L_0$.]
  In the Coulomb $S^3$ construction, the node Laplacian $L_0 = D - A$
  on the Fock graph computes the physical energy spectrum directly:
  the eigenvalues of $L_0$ are $\lambda_n = -(n^2 - 1)$, which map
  to $E_n = -1/(2n^2)$ via the Fock projection.  The graph
  Laplacian \emph{is} the kinetic operator (under the $\kappa$
  calibration).  In the HO $S^5$ Bargmann--Segal construction, $L_0$
  encodes \emph{dipole transition amplitudes}, not the spectrum: the
  spectrum is read off from the holomorphic degree $N$ of each
  node, not from any eigenvalue of $L_0$.  This is the first layer
  of the asymmetry.
\item[Layer 2: Calibration $\pi$.]  Coulomb requires it (yes,
  Section~\ref{sec:transcendental}); HO does not (no, by Paper~24
  $\pi$-free certificate).  This is the transcendental layer.
\item[Layer 3: Wilson physical content of the gauge sector.]  The
  Sprint~5 result~\cite{paper25} on $S^5$ gauge structure showed
  that the abelian U(1) Wilson--Hodge structure transfers verbatim
  to the Bargmann--Segal graph (signed incidence, edge Laplacian
  $L_1 = B^T B$, Hodge decomposition, all hold), but the
  non-abelian SU(3) extension is not natural: transitions between
  $(N, 0)$ and $(N+1, 0)$ irreps are Clebsch--Gordan intertwiners,
  not group elements, so the Wilson lattice gauge requirement of a
  fixed group on every link fails.  By contrast, the Coulomb $S^3$
  construction supports the full SU(2) Wilson lattice gauge of
  Paper~30.  The combinatorial vocabulary of Wilson--Hodge transfers
  to both spheres, but only Coulomb $S^3$ supports a non-trivial
  matter-propagator role for the node Laplacian.
\end{description}

\subsection{The three-layer table}

\begin{table}[h]
\caption{Three-layer Coulomb/HO partition (Sprint~5
sharpening~\cite{paper25}).  $\checkmark$~means the feature is
present; ``--'' means it is absent or not natural.}
\label{tab:three_layers}
\begin{ruledtabular}
\begin{tabular}{lcc}
Feature & Coulomb $S^3$ & HO $S^5$ \\
\hline
$L_0$ computes spectrum & $\checkmark$ & -- \\
Calibration $\pi$ & $\checkmark$ & -- \\
Non-abelian Wilson gauge & $\checkmark$ (SU(2)) & -- \\
\hline
Abelian U(1) Wilson--Hodge & $\checkmark$ & $\checkmark$ \\
Angular sparsity & $\checkmark$ & $\checkmark$ \\
Block-diagonal ERI & $\checkmark$ & $\checkmark$ \\
\end{tabular}
\end{ruledtabular}
\end{table}

The features above the divide in Table~\ref{tab:three_layers} are
the three potential-specific layers; the features below are the
universal-sector features that hold on either sphere.

\subsection{Why three and not two}

The third layer is structurally distinct from the first two.  Layer~1
(spectrum-computing) is a question about \emph{what $L_0$ does}; it
is determined by the projection rule (Fock vs.\ Bargmann--Segal).
Layer~2 (calibration $\pi$) is a question about \emph{which scalars
appear}; it is determined by the operator order and projection type
(Section~\ref{sec:transcendental}).  Layer~3 (Wilson physical content)
is a question about \emph{whether the gauge connection on edges is
group-valued}; it is determined by the representation theory of the
ambient symmetry group (whether transitions between irreps are
$\mathrm{SO}(4)$-equivariant or $\mathrm{SU}(3)$-CG-intertwined).
Three independent structural questions, three independent
discriminators.  None reduces to either of the others.

\medskip\noindent\textbf{Predictive consequence.}  For any candidate
new central potential, the three-layer table provides a checklist
for what features of the GeoVac framework can be expected to
transfer.  A new potential matching the Coulomb pattern (e.g., a
hidden $\mathrm{SO}(4)$-like symmetry, a stereographic-style
projection, integer spectrum) inherits all three Coulomb-specific
layers.  A potential matching the HO pattern (e.g., a hidden
SU($n$)-like symmetry, a complex-analytic projection, linear
spectrum) inherits all three HO-specific layers.  Any other potential
inherits the universal sector but neither set of specific layers,
and must be analyzed system by system.

%=================================================================
\section{Implications}
\label{sec:implications}
%=================================================================

The partition framework simplifies the question of how to extend
GeoVac to a new system.  We list the practical implications.

\subsection{What transfers for free}

Any new central-potential extension of GeoVac inherits the universal
sector without re-derivation:

\begin{itemize}
\item Angular sparsity bounds at every $l_{\max}$
  (Eq.~\eqref{eq:eri_density})~\cite{paper22}.  At $l_{\max} = 3$,
  the new system's two-body Hamiltonian is at least $98.56\%$
  sparse before any potential-specific feature is invoked.
\item Composed-block $O(Q^{2.5})$ Pauli scaling~\cite{paper14}.  The
  composed architecture's structural sparsity advantage over
  Gaussian baselines (51$\times$ to 1712$\times$ across the GeoVac
  molecule library) carries over to any spherical-fermion system.
\item Selection rules for transitions, oscillator strengths,
  multipole moments, and the Wigner--Eckart content of any
  spherical tensor operator.
\item Block-diagonal structure of two-body interaction tensors in
  the composed architecture~\cite{paper17}.
\item $jj$-coupled spinor angular sparsity for relativistic
  extensions~\cite{paper14}~\S V.
\end{itemize}

These features hold by virtue of the algebra $\mathcal{A}$ and the
spherical-harmonic basis decomposition; they are independent of the
choice of $D$ and therefore independent of the radial potential.

\subsection{What requires per-system rederivation}

Every result that depends on the explicit form of $D$ must be
re-derived for the new system:

\begin{itemize}
\item The spectrum $\lambda_n$ of the graph Laplacian (no analog of
  $\lambda_n = -(n^2 - 1)$ unless the system has $\mathrm{SO}(4)$
  symmetry).
\item The kinetic calibration $\kappa$ (no analog of $-1/16$
  unless the projection is Fock-style).
\item The transcendental content of spectral
  invariants~\cite{paper18,paper28}.  The Coulomb spectral zeta
  function lives in the $\pi$-polynomial ring (Paper~28 T9
  theorem); other systems may live in different rings (e.g., the HO
  Bargmann--Segal lattice lives in $\mathbb{Q}$).
\item The fine-structure-constant conjecture
  $K = \pi(B + F - \Delta)$~\cite{paper2_alpha}.  This is a
  Coulomb-specific observation tied to the Hopf-bundle measure
  factor $\pi$ on $S^3$ and has no analog on $S^5$ or other
  sphere-quotients.  Specifically, the Sprint~5 investigation of
  $S^5$ analogs found that no $S^5$ candidate combination matches
  $\alpha^{-1}$ at comparable precision, and the structural mechanism
  is absent (the F-analog candidate
  $(1/2)[\zeta_R(4) + 3\zeta_R(5) + 2\zeta_R(6)]$ is mixed even/odd
  zeta with no $\pi^2/6$ factor).
\item Specific selection rules requiring the explicit form of
  the Dirac operator (e.g., the seven of eight QED selection rules
  recovered by promoting the photon to vector $(q, m_q)$ quantum
  numbers, Paper~28's vector-photon QED
  section~\cite{paper28}, is a Coulomb-$S^3$ result; the $S^5$
  analog would have to be worked out separately).
\end{itemize}

\subsection{What is currently unknown}

Several extensions are not yet placed in the partition framework:

\begin{itemize}
\item Non-central potentials (multiparticle interactions with
  explicit angular dependence).  The angular sparsity theorem
  assumes the two-body interaction is rotationally invariant; for
  non-central interactions the universal sector may shrink.
\item Rank-2 tensor interactions (e.g., Breit retarded interactions,
  spin-orbit, dipole--dipole).  These respect the Gaunt selection
  rules for their angular content but require additional reduced
  matrix elements; whether the universal sector remains the same
  is established in Paper~22~\S V (spinor extension) for the
  $jj$-coupled relativistic case but not for general rank-2 tensor
  interactions.
\item Yang--Mills sectors with matter content (Paper~30~\cite{paper30}
  is purely gauge; matter-coupled extensions have not been
  classified).  The non-abelian Wilson construction is a
  Coulomb-specific Layer-3 feature; whether it admits a HO-side
  analog (perhaps via the $\mathrm{SU}(3)$ symmetry of the
  isotropic oscillator) is an open question.
\end{itemize}

%=================================================================
\section{Open Questions}
\label{sec:open}
%=================================================================

The partition framework as stated is descriptive: it sorts existing
results into universal and potential-specific buckets and predicts
which results carry over to new systems.  It is not yet predictive
in the deeper sense of establishing the partition as a theorem.
Three open questions follow.

\subsection{Is the partition exhaustive?}

The framework as stated identifies two known potential-specific
projections (Fock for Coulomb, Bargmann--Segal for HO) and asserts
that they sit on opposite sides of an operator-order divide.  Are
there other rigidity theorems we have missed?  Specifically:

\begin{itemize}
\item The Morse oscillator $V(r) = D_e(1 - e^{-a(r - r_e)})^2$ has
  an exact analytic spectrum with $\mathrm{SU}(2)$ dynamical
  symmetry but is one-dimensional; whether a 3D radial Morse-like
  potential admits a sphere-quotient discretization with a rigidity
  theorem analogous to Theorems~\ref{thm:fock_rigidity}
  and~\ref{thm:bs_rigidity} is unknown.
\item The P\"oschl--Teller potential and other shape-invariant
  potentials~\cite{infeld_hull1951,cooper_khare1995} have algebraic
  ladder structure analogous to the HO; whether any of them admit
  a Bargmann-style discretization on a non-$S^5$ Hardy space is
  unknown.
\item The Coulomb-Sturmian basis~\cite{paper8_9} is closer to the
  $S^3$ Coulomb construction but uses a different boundary
  condition; it has its own selection rules and transcendental
  structure that have not been placed in the partition framework.
\end{itemize}

A complete classification of central potentials by sphere-quotient
projection type would put the partition framework on firm
mathematical footing.

\subsection{What is the universal/specific partition for non-central
interactions?}

The universal sector as stated assumes rotationally invariant
two-body interactions.  For non-central interactions (gauge
couplings, spin--orbit, multipole moments coupled to external
fields), the angular structure of the interaction itself enters the
selection rules.  In some cases (e.g., scalar QED on $S^3$,
Paper~28~\cite{paper28}) the non-central content is captured by a
modified vertex coupling tensor that respects a residual
$\mathrm{SO}(4)$ structure.  In other cases (Breit retarded
interactions in Paper~14~\S V~\cite{paper14}) it requires an
explicit reduced matrix element that goes beyond the Gaunt algebra.
Whether there is a general universal/specific partition for
non-central interactions, analogous to the rotationally invariant
case, is not yet established.

\subsection{Does the partition predict $\alpha$?}

The fine-structure constant conjecture
$K = \pi(B + F - \Delta) = \alpha^{-1}$~\cite{paper2_alpha} sits in
the potential-specific Coulomb sector by the partition framework.
The headline observation matches CODATA at
$8.8 \times 10^{-8}$~\cite{paper2_alpha}, but the combination rule
itself remains a numerical observation without first-principles
derivation~\cite{curve_fit_audit}.  The Sprint~5 investigation of
$S^5$ analogs~\cite{paper25} confirmed that no analogous combination
exists on the harmonic-oscillator side, consistent with the
partition placing $K$ in the Coulomb-specific sector.  Whether the
partition framework can be sharpened to \emph{predict} the
combination rule, or whether $K$ is genuinely a structural
coincidence among three independent Coulomb-side
ingredients~\cite{paper2_alpha}, is the most consequential open
question in the framework.

%=================================================================
\section{Conclusion}
\label{sec:conclusion}
%=================================================================

The GeoVac framework is a discrete spectral triple in the
graph-spectral-triple lineage of Marcolli and van
Suijlekom~\cite{marcolli_vs2014}, and its results decompose cleanly
along the standard $\mathcal{A}$ versus $D$ split.  The universal
sector --- containing angular sparsity bounds, composed-block Pauli
scaling, selection rules, and the $jj$-coupled spinor extension ---
depends only on the algebra $\mathcal{A}$ and the angular structure
of its basis.  These results transfer to any spherical-fermion system
without re-derivation.  The potential-specific sector --- containing
the spectrum, the calibration $\kappa$, the transcendental content,
the rigidity theorems, and the $\alpha$-conjecture --- depends on
the explicit form of $D$.  These results must be re-derived for each
new system.

The two known potential-specific rigidity theorems (Fock for
Coulomb, Bargmann--Segal for HO) sit on opposite sides of an
operator-order divide.  Coulomb is a second-order Riemannian
construction with a nonlinear projection that introduces calibration
$\pi$; HO is a first-order complex construction with a linear-affine
projection that is bit-exactly $\pi$-free.  The Sprint~5 sharpening
shows the asymmetry has three layers, not two: spectrum-computing
role of $L_0$, calibration $\pi$, and Wilson physical content of the
non-abelian gauge sector.  All three layers stratify along the same
operator-order divide.

The framing offered by this paper has no new theoretical content; it
is a structural reorganization of existing results in
spectral-triple language.  Its value is operational: it tells future
extensions of the framework which results may be quoted from prior
papers and which require system-specific computation, and it places
the framework in a lineage where the relevant tools (graph spectral
triples, gauge networks, spectral truncations) already exist in the
mathematical literature.

\medskip\noindent\textbf{Acknowledgement of source material.}  Every
theorem cited in this paper is established in
Papers~22~\cite{paper22}, 23~\cite{paper23}, 24~\cite{paper24},
or 25~\cite{paper25}.  The contribution here is the explicit
$\mathcal{A}$/$D$ partition framing and the predictive checklist
following from the three-layer table
(Table~\ref{tab:three_layers}).  Numerical verifications of the
universal sector are in Paper~22~\cite{paper22} (five potentials at
$l_{\max} = 2$); numerical verifications of the $\pi$-free
HO-specific sector are in Paper~24~\cite{paper24}
(Bargmann--Segal lattice at $N_{\max} = 5$, 56 nodes, 165 edges,
zero irrationals); numerical verification of the $S^5$ gauge
structure non-naturalness is in the Sprint~5 record~\cite{paper25}.

%=================================================================
\begin{thebibliography}{99}
%=================================================================

\bibitem{connes_book1994}
A.~Connes,
\textit{Noncommutative Geometry}
(Academic Press, San Diego, 1994).

\bibitem{chamseddine_connes1997}
A.~H.~Chamseddine and A.~Connes,
``The Spectral Action Principle,''
Commun.\ Math.\ Phys.\ \textbf{186}, 731 (1997).

\bibitem{marcolli_vs2014}
M.~Marcolli and W.~D.~van Suijlekom,
``Gauge networks in noncommutative geometry,''
J.\ Geom.\ Phys.\ \textbf{75}, 71 (2014); arXiv:1301.3480.

\bibitem{perez_sanchez2024}
C.~I.~P\'erez-S\'anchez,
``On the gauge networks of Marcolli--van Suijlekom: Yang--Mills
without Higgs,'' arXiv:2401.03705 (2024); see also arXiv:2508.17338
for the corrected continuum limit analysis.

\bibitem{lichnerowicz1963}
A.~Lichnerowicz,
``Spineurs harmoniques,''
C.~R.\ Acad.\ Sci.\ Paris \textbf{257}, 7 (1963).

\bibitem{bander_itzykson1966}
M.~Bander and C.~Itzykson,
``Group Theory and the Hydrogen Atom,''
Rev.\ Mod.\ Phys.\ \textbf{38}, 330 and 346 (1966).

\bibitem{fock1935}
V.~Fock,
``Zur Theorie des Wasserstoffatoms,''
Z.\ Phys.\ \textbf{98}, 145 (1935).

\bibitem{bargmann1961}
V.~Bargmann,
``On a Hilbert space of analytic functions and an associated
integral transform, Part I,''
Commun.\ Pure Appl.\ Math.\ \textbf{14}, 187 (1961).

\bibitem{hall1994}
B.~C.~Hall,
``The Segal--Bargmann coherent state transform for compact Lie
groups,''
J.\ Funct.\ Anal.\ \textbf{122}, 103 (1994).

\bibitem{jauchhill1940}
J.~M.~Jauch and E.~L.~Hill,
``On the Problem of Degeneracy in Quantum Mechanics,''
Phys.\ Rev.\ \textbf{57}, 641 (1940).

\bibitem{borelweil}
A.~Borel and A.~Weil,
``Repr\'esentations lin\'eaires et espaces homog\`enes
k\"ahleriens des groupes de Lie compacts,''
S\'eminaire Bourbaki, expos\'e 100 (1954).

\bibitem{Dyall}
K.~G.~Dyall and K.~F\ae gri,
\textit{Introduction to Relativistic Quantum Chemistry}
(Oxford University Press, Oxford, 2007).

\bibitem{infeld_hull1951}
L.~Infeld and T.~E.~Hull,
``The Factorization Method,''
Rev.\ Mod.\ Phys.\ \textbf{23}, 21 (1951).

\bibitem{cooper_khare1995}
F.~Cooper, A.~Khare, and U.~Sukhatme,
``Supersymmetry and Quantum Mechanics,''
Phys.\ Rep.\ \textbf{251}, 267 (1995).

\bibitem{paper0}
J.~Loutey,
``Geometric Packing and the Universal Constant $K = -1/16$,''
GeoVac Paper~0 (2026).

\bibitem{paper1}
J.~Loutey,
``Spectral Graph Theory and the GeoVac Lattice,''
GeoVac Paper~1 (2026).

\bibitem{paper2_alpha}
J.~Loutey,
``The Fine Structure Constant from Spectral Geometry of the Hopf
Fibration,''
GeoVac Paper~2 (2026). [Observation; combination rule conjectural.]

\bibitem{paper7}
J.~Loutey,
``The Dimensionless Vacuum: Spectral Graph Theory and Conformal
Equivalence on $S^3$,''
GeoVac Paper~7 (2026).

\bibitem{paper8_9}
J.~Loutey,
``Bond Sphere Geometry and the Sturmian Structural Theorem,''
GeoVac Papers~8--9 (2026).

\bibitem{paper13}
J.~Loutey,
``Hyperspherical Lattice for Multi-Electron Atoms,''
GeoVac Paper~13 (2026).

\bibitem{paper11}
J.~Loutey,
``Prolate Spheroidal Lattice for Diatomic Molecules,''
GeoVac Paper~11 (2026).

\bibitem{paper14}
J.~Loutey,
``Structurally Sparse Qubit Hamiltonians from Spectral Graph
Theory,''
GeoVac Paper~14 (2026).

\bibitem{paper15}
J.~Loutey,
``Level 4 Mol-Frame Hyperspherical Geometry for Diatomic
Molecules,''
GeoVac Paper~15 (2026).

\bibitem{paper17}
J.~Loutey,
``Composed Natural Geometries for Core-Valence Molecular
Electronic Structure,''
GeoVac Paper~17 (2026).

\bibitem{paper18}
J.~Loutey,
``Exchange Constants and the Transcendental Taxonomy of
Observables,''
GeoVac Paper~18 (2026).

\bibitem{paper22}
J.~Loutey,
``Potential-Independent Angular Sparsity for Qubit Hamiltonians in
Spherical Fermion Systems,''
GeoVac Paper~22 (2026).

\bibitem{paper23}
J.~Loutey,
``Nuclear Shell Model Qubit Hamiltonians and the Fock Projection
Rigidity Theorem,''
GeoVac Paper~23 (2026).

\bibitem{paper24}
J.~Loutey,
``The Bargmann--Segal Lattice: $\pi$-Free Discretization of the
Harmonic Oscillator on $S^5$,''
GeoVac Paper~24 (2026).

\bibitem{paper25}
J.~Loutey,
``The Hopf Graph as a Lattice Gauge Structure: A Framework
Observation,''
GeoVac Paper~25 (2026).

\bibitem{paper28}
J.~Loutey,
``QED on $S^3$: T9 Theorem, Selection Rules, and the Graph-Native
Formulation,''
GeoVac Paper~28 (2026).

\bibitem{paper30}
J.~Loutey,
``SU(2) Wilson Lattice Gauge on the GeoVac Hopf Graph,''
GeoVac Paper~30 (2026).

\bibitem{curve_fit_audit}
J.~Loutey,
\textit{Curve-fit / pattern-match audit memo},
GeoVac internal document, \texttt{docs/curve\_fit\_audit\_memo.md}
(2026-05-02).

\end{thebibliography}

\end{document}
